ACTS has the potential to grow in multiple directions in the future. It provides a solid foundation for tracking in high-energy physics experiments and can be adapted to meet the needs of different experiments.

One of the areas for expansion is providing tools and algorithms for telescope detector geometries and low-momentum tracking. The current implementation of ACTS is primarily focused on collider experiments, like ATLAS at the LHC, with a lower momentum bound of about 1 GeV. Tracking in the low-momentum regime will require refined material interaction modeling to provide optimal resolutions and the development of new algorithms to fit CPU budget requirements.

The general-purpose track finding algorithm provided in ACTS is a promising starting point for future developments. It has been shown to be efficient in terms of computation and physics performance for the ODD and the ATLAS ITk. However, there are still inefficiencies that seem to be related to the handling of material interactions. The ACTS community is working on analyzing the material description and its impact on track finding performance. At the same time, the track finding algorithm is not yet optimized for electrons that undergo bremsstrahlung. Future developments in track finding in ACTS will focus on this aspect, which is crucial for the performance of ITk track reconstruction and other experiments.

During optimization of the track finding algorithm, it was found that the main track seeding algorithm provided by ACTS is a performance bottleneck. Since it only considers three detector hits, the fake rate is significant in high-pileup conditions, like those expected with the HL-LHC. The ACTS community is working on implementing new seeding algorithms and improving the existing one.

4D tracking is a promising direction for the future of tracking in high-energy physics experiments. The integration of time information into reconstruction has the potential to significantly improve physics performance by reducing the impact of pileup. ACTS provides a solid foundation for 4D tracking, and many algorithms have been adapted to work with time information. Future developments could include adapting the track seeding algorithm to work with time information, which might be crucial for future experiments facing increasing pileup. The developments in 4D vertexing in ACTS already show improvements for pileup mitigation, which are promising results and are about to be published in the near future.

The ACTS community is continuously growing, and new experiments are looking into adopting ACTS for their track reconstruction. To guarantee the success of ACTS in the future, it is important to uphold and strengthen the experiment independence of ACTS. This could include experiment-specific track parameterizations and units in the future. The modularity of ACTS has already been shown to be beneficial for developing new algorithms and adapting them to different experimental setups.
